\section{Relation-index scan: details and sensitivity}
\label{app:stage1}

\paragraph{Run.}
The scan (\texttt{stage1\_v4\_calibrated}) covered all 1,024 heads on 534 prompts
(89 discovery families $\times$ base/corrupted/reorder $\times$ two orders, four
demonstrations each). Preflight checks reproduced the behavioural candidate scores
and showed zero drift under inert hooks and self-patching; attention capture had a
maximum vocabulary-logit drift of 0.039 on one prompt. The test-only extension
(\texttt{ri\_test\_v2}) scored all 46,634 saved test-block events; 237,685
replication checks of saved scores passed. Independent recomputation reproduced all
reported means, contrasts and support counts.

\paragraph{Method details.}
\emph{Event score.} For a qualifying $(T,j)$ with anchor token $t_o$, the score is
$a^{h,j}_T=q^{h,j}_{t_o}/\sum_{k\le j}q^{h,j}_{t_k}$ with
$q^{h,j}_t=\max(0,p^{h,j}_t-\bar p^{h,j})$ and
$p^{h,j}=\mathrm{softmax}(x_jW^h_{OV}W_U)$ restricted to the distinct token ids
visible at $j$ (Appendix~\ref{app:ri}); anchor-token collisions between target and
control names are treated as ties. \emph{Aggregation.} Scores are averaged within
family and then with equal weight across the families that have at least one
scored event for that head; families without events do not enter the mean.
\emph{Name controls.} The other test-fact mothers whose full name is visible at
$j$; token length need not match. \emph{Word controls.} Up to three word types
sampled from the visible non-name tokens of the prompt (mostly template words);
they provide a weaker specificity check. All reported means and counts use the
name-control-eligible subset of events; events without an available control are
excluded from the gap statistics but retained in the pooled index.
\emph{Calibration population.} Test-block events in which both the target and a
control mention are fully visible, of equal token length and with distinct first
tokens; 19,901 of 278,134 QK events qualify. Heads need at least 50 such events
across ten families (80 heads). The null swaps target and control labels
consistently within each family (100{,}000 replicates) and assumes the two labels
are exchangeable under the null; heads are tested with Holm correction over 1,024.
\emph{Selection population.} All name-control-eligible events; a ranking requires
at least 20 events in ten families. The 50-event threshold of the calibration and
the 20-event threshold of the selection are different requirements for different
purposes and yield different head populations.

\paragraph{Why the pooled index is dominated by demonstrations.}
Under the original pooled index and the descriptive 99.9th-percentile rule the two
top heads are L3H11 and L9H22. Table~\ref{tab:pooled} splits their events by block:
demonstrations contribute over 99\% of each head's summed target score, and neither
head has a single event from a test position attending to a demonstration source.
Inspection of the events shows why: all 86 passes of L9H22 attend from a position to
itself, and all 149 passes of L3H11 attend one token back, from a period or newline
to the end of the preceding name. In both cases the OV projection is computed from
the same current-token embedding for every event, so the reported index varies only
with which target token is selected and how the visible context normalizes it. A
high pooled index can thus be obtained from local attention with a static token
preference, without any event that demonstrates retrieval of a test relation. This
does not show that these heads take no part in relational retrieval; it shows that
the pooled score does not require it, which motivates the test-block restriction of
\S\ref{sec:method-ri}.

\begin{table}[h]
\centering\scriptsize
\setlength{\tabcolsep}{3pt}
\begin{tabular}{@{}lrrrr@{}}
\toprule
Head & Pooled RI & QK freq.\ (\%) & Demo: $n$ / RI & Test: $n$ / RI \\
\midrule
L3H11 & 0.243 & 0.0097 & 144 / 0.250 & 5 / 0.038 \\
L9H22 & 0.374 & 0.0056 & 72 / 0.445 & 14 / 0.007 \\
\bottomrule
\end{tabular}
\caption{The two heads selected by the pooled index and the descriptive
99.9th-percentile rule, split by whether the annotated fact and the current position
lie in a demonstration or in the test block. Frequency is the share of eligible
fact--position evaluations that pass the QK gate.}
\label{tab:pooled}
\end{table}

\paragraph{Matched-target calibration.}
The matched subset retained 19,901 of 278,134 QK events; 80 heads had at least 50
matched events across ten families. Under 100{,}000 family-consistent target/control
swaps, no head passed Holm correction at $\alpha{=}0.05$, over 1,024 or over the 80
supported heads. The three smallest raw $p$ values are L16H4 ($p{=}0.006$, target
0.0147 vs.\ control 0.0093, 87 events, 40 families), L9H18 ($p{=}0.023$) and L13H13
($p{=}0.025$). Neither pooled-index head had a matched event.

\paragraph{Test-only candidates.}
Table~\ref{tab:ricands} lists the leading heads under the test-only index with name
controls (all test facts, all test positions). The frozen rule (\S\ref{sec:method-ri})
selects 59 heads in total, 36 of which enter at least one name-gap ranking. Two
patterns recur among high scorers. Some heads owe much of their score to events in
which the current token is part of the target name itself (L23H10: 31 of 60
comparisons; L15H18: 104 of 110); such events can be scored without any retrieval
across positions, so they do not demonstrate it, although they do not exclude it.
Excluding those events reduces but does not remove L23H10's gap. Others depend on one current token (for L17H5 the
token \emph{mother} contributes 66\% of the family-weighted target score) or on the
reorder variant (299 of L9H16's 343 comparisons). Paired base/corrupted events
that follow the changed fact are rare: L11H4 has 21 jointly passing pairs in 18
families, with a positive name gap in both versions for only five families.

\begin{table}[h]
\centering\scriptsize
\setlength{\tabcolsep}{3pt}
\begin{tabular}{@{}llrrrrr@{}}
\toprule
Head & Anchor & Events & Fam. & Target & Control & Gap \\
\midrule
L11H4  & first & 94  & 39 & 0.0147 & 0.0081 & $+0.0065$ \\
L17H5  & first & 65  & 34 & 0.0229 & 0.0168 & $+0.0061$ \\
L9H16  & last  & 343 & 87 & 0.0209 & 0.0135 & $+0.0074$ \\
L12H2  & first & 52  & 32 & 0.0201 & 0.0129 & $+0.0072$ \\
L25H18 & last  & 46  & 17 & 0.0287 & 0.0156 & $+0.0131$ \\
L23H10 & first & 60  & 24 & 0.0283 & 0.0068 & $+0.0215$ \\
\bottomrule
\end{tabular}
\caption{Leading test-only candidates (all test facts and positions). Events are
target/control comparisons; family-weighted means over the name-control-eligible
subset.}
\label{tab:ricands}
\end{table}

\paragraph{The answer position and the raw-embedding invariance.}
At the final colon, 20 heads have qualifying events (748 in total, 744 to test
sources); L16H21 and L16H1 attend to the query fact in over 220 prompts each (41\%
of prompts). Their name gaps at that position are near zero or negative
(Table~\ref{tab:final}). Part of the reason is structural. The current token at the answer position is
always the colon, so for a fixed head the raw-embedding projection vector is the same
in every prompt, and the normalized score is the same whenever the set of visible
token ids is the same; only the attention condition can differ between such prompts.
In all 49 jointly passing base/corrupted pairs of L16H21 the sign of the
old-versus-new-mother contrast flips exactly when the target label flips, at both
anchors: once the head attends, the scoring component cannot follow a reassignment
of mothers at this position. The same holds at any earlier position whose current
token and visible token set are unchanged. This is a limit of the score, not a
statement about what the head's actual output carries.

\begin{table}[h]
\centering\small
\begin{tabular}{@{}lrrrrr@{}}
\toprule
Head & Query events & Fam. & QK (\%) & First gap & Last gap \\
\midrule
L16H21 & 222 & 77 & 41.6 & $-0.0004$ & $+0.0034$ \\
L16H1  & 223 & 76 & 41.8 & $-0.0025$ & $-0.0034$ \\
L16H24 & 68  & 35 & 12.7 & $-0.0028$ & $+0.0036$ \\
L16H4  & 63  & 32 & 11.8 & $+0.0004$ & $-0.0072$ \\
\bottomrule
\end{tabular}
\caption{Heads whose attention at the answer position selects the query fact
(denominator 534 prompts), with their name gaps at the two anchors.}
\label{tab:final}
\end{table}

\begin{figure}[h]
\centering
\includegraphics[width=\columnwidth]{figE_ri_distribution}
\caption{The pooled index over heads. (a) Distribution of the first-anchor index for
the 342 heads with at least one qualifying event (log count); the two heads selected by
the percentile rule are marked. (b) First- against last-anchor index for all heads;
the anchor changes which heads lead (L29H14 tops the last-anchor ranking).}
\label{fig:ri-dist}
\end{figure}

\paragraph{Dominance margin.}
Over all 1,455,940 eligible dominance ratios (event-weighted, demonstrations
included), 80.9\% are at or below the published $\tau{=}2.2$ and the 95th percentile
is 4.48. Refiltering the saved passes at 4.48 keeps 72,797 events; the two pooled
heads keep 30 and 49 passes respectively. The anchor also matters: under the last
token, L29H14 tops the conditional-strength ranking (Figure~\ref{fig:ri-dist}b). All such variants are reported
as sensitivity analyses; the audit itself keeps the published parameters.
